\section{Davis--Kahan result inventory}
\label{app:dk-inventory}

\Cref{tab:dk-full-inventory} records the maintained result-level denominator.
Every row is currently verified in the ordinary Lean build and has accepted
semantic/source-correspondence review. The disposition column distinguishes a
proof of the source claim from the checked refutation of Proposition~4.4.

\begin{longtable}{>{\raggedright\arraybackslash}p{0.16\linewidth}>{\raggedright\arraybackslash}p{0.55\linewidth}>{\raggedright\arraybackslash}p{0.18\linewidth}}
\caption{The \DKResultCount{} Davis--Kahan result-level targets.}
\label{tab:dk-full-inventory}\\
\toprule
ID & Result & Disposition \\
\midrule
\endfirsthead
\toprule
ID & Result & Disposition \\
\midrule
\endhead
\\path{S2-sin-theta} & Single-angle sine theorem & proved at source scope \\
\\path{S2-tan-theta} & Single-angle tangent theorem & proved at source scope \\
\\path{S2-sin-two-theta} & Double-angle sine theorem & proved at source scope \\
\\path{S2-tan-two-theta} & Double-angle tangent theorem & proved at source scope \\
\\path{DK-3.1-prop} & Acute direct rotation existence and uniqueness & proved at source scope \\
\\path{DK-3.2-prop} & Nonacute existence criterion & proved at source scope \\
\\path{DK-3.3-prop} & Principal square-root characterization & proved at source scope \\
\\path{DK-3.4-prop} & Square as a direct rotation & proved at source scope \\
\\path{DK-3.1-thm} & Classification of pairs of subspaces & proved at source scope \\
\\path{DK-3.1-cor} & Compact classification by angle eigenvalues & proved at source scope \\
\\path{DK-3.5-prop} & Angle commutation and eigenspace geometry & proved at source scope \\
\\path{DK-3.2-cor} & Reversal symmetry & proved at source scope \\
\\path{DK-4.1-prop} & Pointwise and singular-value extremality of the direct rotation & proved at source scope \\
\\path{DK-4.1-cor} & Unitarily invariant norm minimality of restricted displacement & proved at source scope \\
\\path{DK-4.2-prop} & Basis-angle square-sum extremality & proved at source scope \\
\\path{DK-4.3-prop} & Squared-displacement unitarily invariant norm minimality & proved at source scope \\
\\path{DK-4.4-prop} & Full-displacement claim of Proposition~4.4 & refuted as printed; repaired separately \\
\\path{DK-5.1-thm} & Banach-space Sylvester lower bound & proved at source scope \\
\\path{DK-5.2-thm} & Semibounded self-adjoint Sylvester theorem & proved at source scope \\
\\path{DK-5.1-lem} & Strong-cutoff convergence of singular values & proved at source scope \\
\\path{DK-6.1-lem} & Direct-sum norm comparison and converse & proved at source scope \\
\\path{DK-6.2-lem} & Reflection-pinch contraction & proved at source scope \\
\\path{DK-6.1-prop} & Sine proof, ambient limitation, and symmetric sine theorem & proved at source scope \\
\\path{DK-6.1-thm} & Generalized sine theorem & proved at source scope \\
\\path{DK-6.2-thm} & Pairwise-gap square-norm sine theorem & proved at source scope \\
\\path{DK-6.3-thm} & Generalized tangent theorem and supporting machinery & proved at source scope \\
\\path{DK-6.3-lem} & Finite-rank near-maximizer leakage estimate & proved at source scope \\
\\path{DK-8.1-thm} & Branch selection and spectral repulsion & proved at source scope \\
\\path{DK-8.2-thm} & Smallness selects the acute branch & proved at source scope \\
\bottomrule
\end{longtable}

The result table is keyed to source results.  One source result can require many
Lean declarations, while a single reusable Lean theorem can support more than
one source-facing wrapper.  The denominator therefore remains stable under
refactoring of implementation modules.

\section{Source-correspondence review}
\label{app:source-review}

A result is accepted only after two distinct checks.

\paragraph{Mechanical verification.}
The declaration must resolve in the ordinary project build, and every theorem
on which it depends must be accepted by Lean's kernel.  This establishes the
formal proposition exactly as encoded.

\paragraph{Source comparison.}
The encoded proposition is compared with the cited source location.  The review
checks the mathematical objects, hypotheses, conclusion, scalar field,
dimension, norm class, boundary conventions, and operator-domain assumptions.
For a nonliteral representation, the review also records the theorem or
mathematical argument that justifies the correspondence.  Equation
\cref{eq:modulus-sine} is the main example: the rectangular map in the Lean
sine-theta theorem represents the source's positive sine operator only after a
modulus/polar-decomposition argument.

This separation prevents theorem names from serving as evidence of fidelity.
For example, the proved declaration
\texttt{sinTheta\_unbounded\_intervalExterior\_characterizedWitness\_rclike}
remains a useful finite-gap theorem, but it is not the selected source endpoint
because it does not include the half-infinite gap cases.  Conversely,
Proposition~4.4 is complete in the inventory even though its printed assertion
has no proof: the completed formal treatment consists of the transcribed claim
as a proposition, a checked counterexample/refutation, a diagnosis of the
source proof, and repaired theorems.

\section{YWS source-audit scope}
\label{app:yws-coverage}

The YWS audit uses a different denominator because the journal paper does not
claim a line-by-line formalization of every mathematical statement in that
article.  The maintained census has \YWSCensusEntryCount{} rows:
\YWSCensusPaperSourceCount{} rows associated with printed source material and
\YWSCensusExplicitAdditionCount{} explicitly marked additions.  All
\YWSCensusVerifiedCount{} rows currently resolve to declarations in the normal
build.

The source rows include the conclusions of Theorems~1--3 and the equations,
lemmas, examples, and boundary conventions needed to check those conclusions.
Several rows can share one correction.  In particular, the current census has
\YWSCensusCorrectedCount{} rows whose implementation is classified as corrected,
but the manuscript reports two printed YWS defects: the missing square in
equation~(4) and the rank-boundary convention in Theorem~3.

\section{Formalization procedure and AI assistance}
\label{app:formalization-process}

The development used a recurring five-stage workflow.
\begin{enumerate}
  \item \textbf{Gather context.}  Read the source result and surrounding proof,
        identify standing assumptions and notation, and inspect the relevant
        Lean interfaces.
  \item \textbf{Decompose.}  Express the target as mathematical obligations
        small enough to compare individually with existing library facts.
  \item \textbf{Find foundations.}  Search Mathlib, Tau Ceti, the local reusable
        layer, other Lean developments, and the mathematical literature before
        introducing new theory.
  \item \textbf{Formalize.}  Implement definitions and lemmas, compile them with
        Lean, and revise failures until the desired theorem is kernel-checked.
  \item \textbf{Review skeptically.}  Compare the resulting statement with the
        source independently of successful compilation.  Reopen the result if
        the review finds a scope, convention, representation, or hypothesis
        mismatch.
\end{enumerate}
The loop terminates for a source target only when both the formal proof and the
source-correspondence review are accepted, or when the source statement is
shown false and the formal record contains the corresponding refutation.

Large language models were used for proof planning, repository search, Lean
implementation, proof repair, and source comparison.  Human review determined
the intended paper scope, evaluated mathematical correspondence, and accepted
or rejected proposed repairs.  All formal propositions are certified by Lean's
kernel; model outputs supplied proof text and analysis that remained subject to
compilation and review.

\section{Artifact organization}
\label{app:artifact}

The paper-facing Davis--Kahan declarations are under
\path{DavisKahan/Sources/DavisKahan1970}.  Reusable mathematical infrastructure
is primarily under \path{ForTauCeti}; its declarations already use final
\texttt{TauCeti.*} namespaces.  The maintained Davis--Kahan result inventory is
\path{dev/davis-kahan-1970-formalization-result-inventory.json}.  The YWS source
census is maintained separately under the YWS package and is summarized for
this manuscript by the generated census snapshot.  This organization permits
source-audit records and paper-specific counterexamples to remain local to the
paper package while reusable analysis is prepared for upstream integration.
